\documentclass[11pt]{article}

\usepackage[a4paper,margin=1in]{geometry}
\usepackage[T1]{fontenc}
\usepackage[utf8]{inputenc}
\usepackage{lmodern}
\usepackage{microtype}
\usepackage{amsmath,amssymb}
\usepackage{booktabs}
\usepackage{xcolor}
\usepackage{graphicx}
\usepackage[hidelinks]{hyperref}

\newcommand{\sandboxer}{\textsc{Sandboxer}}
\newcommand{\bluephase}{Blue Phase}
\newcommand{\redphase}{Red Phase}

\title{\Large \textbf{Sandboxer: An Experimental Model-versus-Model Benchmark for\\ Offensive and Defensive Cyber Capability in a Simulated Capture-the-Flag Arena}}

\author{The Sandboxer Team\\ \small \texttt{https://github.com/andreamicheli/sandboxer}}
\date{\today}

\begin{document}

\maketitle

\begin{abstract}
\noindent
\sandboxer{} is an experimental benchmark that evaluates the offensive and defensive
cyber capabilities of frontier AI models in a simulated, controlled, and fully
observable capture-the-flag (CTF) arena. Where most existing benchmarks measure
whether a single model can solve a pre-authored security challenge, \sandboxer{}
instead pits two models directly against one another: each model first builds and
defends its own synthetic network service (\bluephase), and then attacks the
service built by its opponent (\redphase), under symmetric, audited budgets.
Vulnerabilities are never planted by the benchmark; they must emerge from the
models' own defensive choices, which is precisely what the arena is designed to
reveal. A trusted Orchestrator and a hybrid deterministic-and-LLM Match Auditor
enforce isolation, symmetry, and validity, while an append-only, hash-chained
telemetry stream acts as the authoritative record from which results, replays,
reports, and broadcast videos are derived. This paper introduces the benchmark,
describes its design, protocol, scoring, evidence, and auditing machinery, and
states the limitations that bound the claims it can support. It is the
introductory announcement of \sandboxer{} ahead of its first public best-of-three
series between two open-weight models, and it is deliberately a design document,
not yet a report of results.
\end{abstract}

\section{Introduction}

Frontier language models are increasingly deployed as autonomous agents that plan,
call tools, and act on external systems over long horizons. As their autonomy
grows, so does the practical importance of understanding what they can do in
security-relevant settings---not only how well they can answer quiz-style
questions about vulnerabilities, but how they behave when they must build a
defense and, symmetrically, when they must break one under time, token, and tool
constraints. Measuring these capabilities rigorously is hard: real systems cannot
be attacked for evaluation purposes, and historical challenge datasets, however
carefully curated, exercise only one side of the coin.

The evaluation community has responded with a family of cybersecurity benchmarks
that score a model's ability to solve curated challenges, among them Cybench
\cite{cybench}, NYU CTF Bench \cite{nyuctf}, CyberGym \cite{cybergym}, and the
EnIGMA agent \cite{enigma}. These are valuable and complementary efforts, but they
share a common structure: a single agent attacks a fixed, pre-authored target,
and the evaluation produces a success rate against that static corpus. They do
not measure how a model defends a system it built itself, nor do they measure how
two models interact adversarially under controlled, symmetric conditions.

\sandboxer{} is designed to fill this gap. Its core unit is a \emph{Match}: two
models are placed in separate, disposable, isolated virtual machines. In the
\bluephase, each model receives the same minimal synthetic service and the same
seeded ``Blue Brief''---a small functional requirement, not an exploit---and is
free to modify and defend its own service. In the \redphase, the two arenas are
connected along narrow, declared paths, and each model attempts to capture a
synthetic flag from the opponent's service while preserving its own. The
vulnerabilities that matter are therefore \emph{emergent}: they arise from the
models' independent defensive choices, rather than being selected by the seed or
pre-baked into the challenge. This makes the benchmark a test of integrated
offensive and defensive behavior rather than a test of challenge-solving.

Three commitments distinguish \sandboxer{} from a demonstration. First,
\emph{observability}: every state transition, tool call, budget observation, and
submission is recorded in an append-only, hash-chained telemetry stream that is
the authoritative source of truth; the video, the report, and the website are all
derived views of that stream. Second, \emph{symmetry}: the two competitors face
identical prompts, tools, budgets, and provider treatment, and any material
asymmetry invalidates the comparison rather than being papered over. Third,
\emph{bounded claims}: \sandboxer{} is an experimental benchmark, and its public
claim is deliberately narrow---the result describes these specific models, in this
specific configuration, and does not measure general cyber capability or safety
against real-world systems.

This paper is the introductory announcement of the benchmark. We first situate it
against related work (Section~\ref{sec:background}), then describe the arena,
the match protocol, scoring, telemetry, and auditing in detail
(Section~\ref{sec:methodology}), and close with the limitations and threats to
validity that bound what the benchmark can claim (Section~\ref{sec:limitations}).

\section{Background and Related Work}
\label{sec:background}

\subsection{The broader evaluation landscape}

The evaluation of frontier models has matured from static question-answering
benchmarks toward \emph{capability} evaluations that ask how much useful work an
agent can complete over an extended horizon. METR proposes measuring AI
performance by the length of tasks an agent can complete, arguing that
task-completion time horizon is a stable and interpretable metric for progress
\cite{metr}. Epoch AI's FrontierMath demonstrates the value of expert-crafted,
unseen, and independently vetted problems for probing capability at the frontier
\cite{frontiermath}. The UK AI Safety Institute's Inspect framework provides an
open, extensible harness for running and logging such evaluations reproducibly
\cite{inspect}. In the software-engineering domain, SWE-bench \cite{swebench} and
InterCode \cite{intercode} established the standard of scoring an agent by
whether its actions, executed against a real environment, actually achieve the
stated objective.

\sandboxer{} inherits from this tradition the discipline of execution-based
scoring, frozen configurations, and versioned evidence, but applies it to an
adversarial, two-sided task rather than a single-agent completion task.

\subsection{Cybersecurity benchmarks}

A distinct line of work evaluates language-model agents on security tasks.
Cybench provides a framework for specifying cybersecurity tasks and scoring
agents on 40 professional-level CTF challenges \cite{cybench}. NYU CTF Bench
curates 200 CTF challenges across six categories and evaluates both black-box and
open-source models \cite{nyuctf}. CyberGym scales real-world evaluation to 1{,}507
historical vulnerabilities across 188 software projects \cite{cybergym}. EnIGMA
shows that interactive agent-computer interfaces substantially improve a model's
ability to solve CTF challenges \cite{enigma}. Together these efforts establish
that (i) language-model agents possess non-trivial offensive capability on
curated challenges, and (ii) the choice of interface, scaffolding, and budget
materially changes measured performance.

What these benchmarks share is a one-sided structure: a single agent attacks a
pre-authored target, and the result is a per-challenge success rate. Three
dimensions are therefore under-explored. First, \emph{defense}: there is no
symmetrically rigorous measure of how well a model constructs and maintains a
defense. Second, \emph{interaction}: real security is a contest, and a model's
behavior against an adaptive opponent is a different quantity from its behavior
against a static target. Third, \emph{emergence}: when vulnerabilities are
pre-authored, the benchmark cannot distinguish the model's skill from the
difficulty of the planted flaw.

\sandboxer{} is positioned at the intersection of these two lines of work. It uses
the disciplined evaluation machinery of the capability-evaluation tradition
(frozen configuration, execution-based scoring, versioned telemetry, an
independent auditor) applied to the two-sided, emergent task structure that
single-agent CTF benchmarks cannot express. It is closer in spirit to a
controlled ``model duel'' than to a challenge corpus, and it is explicitly a
\emph{showcase evaluation}---a first, carefully bounded answer to the question of
how two models compare when they both defend and attack---rather than a claim of
mature scientific-benchmark status comparable to the established institutions
cited above.

\section{Methodology}
\label{sec:methodology}

\subsection{Design principles}

Four principles govern every design decision in \sandboxer{}.

\begin{itemize}
  \item \textbf{Safety first.} The arena is synthetic and self-contained. Targets
        are toy services and ephemeral flags; there is no path to real systems,
        real credentials, or the control plane. Isolation is deny-by-default and
        enforced by deterministic controls that fail closed.
  \item \textbf{Symmetry.} Both competitors receive identical prompts, tools,
        budgets, provider treatment, and opportunity. Any material asymmetry in a
        measured category invalidates the comparison unless a predeclared
        symmetric fallback removes that category from both sides.
  \item \textbf{Observability.} A single, authoritative, hash-chained telemetry
        stream records everything that matters. Nothing is inferred from screen
        recordings; the replay, the report, and the video are all derived from
        the same validated events.
  \item \textbf{Emergent vulnerability.} No exploit is planted. Both models start
        from the same pristine service and the same seeded functional brief; the
        attack surface that matters is whatever each model's own choices open up.
\end{itemize}

\subsection{Arena and containment}

Each Match runs two disposable Runner virtual machines under a trusted
Orchestrator. Runners are non-privileged, capability-free, and read-only at the
root filesystem, with bounded CPU, memory, and process limits. During the
\bluephase{} their networks are fully separated; a shared, deny-by-default arena
network is attached only at the transition to the \redphase, and even then only
the declared opponent service paths are reachable. There is no egress to the
public internet, no host filesystem mounts, and the Orchestrator---which is never
a valid target---cannot be addressed, terminated, or administered from the arena.
Flags, services, and targets are synthetic; nothing outside the arena is a
permitted objective. Every terminal state ends in verified teardown, and any
teardown that cannot be confirmed destroys the validity of the Match.

\subsection{Task family and the Blue Brief}

The task is deliberately minimal so that the interesting variation comes from the
models, not from an elaborate scenario. Both competitors begin from the same
small, pristine, synthetic notes service with deterministic health,
functional-integrity, and flag-capture tests. The service contains no planted
exploit.

At the start of a Match, both competitors receive the same seed-selected
\textbf{Blue Brief}: a minor functional requirement---such as constrained public
export, controlled sharing, or sanitized diagnostics---stated as an outcome, not
an implementation. Each model independently changes and defends its own service
to satisfy the brief. The briefs are versioned, frozen for a series, and
calibrated to be symmetric and of comparable difficulty. Because the brief
specifies \emph{what} to add and not \emph{how} to secure it, the defensive
choices---and therefore the attack surface---emerge from the model. A best-of-three
series uses three distinct briefs selected without replacement from the frozen
catalog.

\subsection{Match protocol}

A Match is a validated state machine controlled by the Orchestrator:

\begin{center}
\begin{minipage}{0.92\textwidth}
\begin{quote}
\small
\texttt{specified $\rightarrow$ provisioned $\rightarrow$ preflight\_passed
$\rightarrow$ blue\_active $\rightarrow$ blue\_validation
$\rightarrow$ intermission\_interview $\rightarrow$ flags\_injected
$\rightarrow$ red\_preflight $\rightarrow$ red\_active $\rightarrow$ finalizing
$\rightarrow$ valid\_result $|$ invalid $|$ aborted $\rightarrow$ teardown\_verified}
\end{quote}
\end{minipage}
\end{center}

Every transition is an authoritative telemetry event with a monotonic timestamp,
a reason, and a responsible component, and a failed prerequisite cannot be
bypassed by advancing the state manually.

\paragraph{Blue Phase.} Both competitors build and defend their own service in
independent concurrent loops. Allowed tools are limited to inspecting, editing,
restarting, and testing the model's \emph{own} service, plus an irreversible
\texttt{finish\_phase}. The phase ends when both competitors finish or exhaust
their Blue budgets. The Auditor then verifies that both services pass the same
minimum health contract; if either fails, the Red Phase does not start and the
Match is invalid rather than an automatic competitive loss.

\paragraph{Intermission Interview.} Between phases, each model receives the same
separate, tool-free interview asking how its defensive phase went, which
weaknesses it perceives, and what attack it expects. The interview uses a small
equal budget, cannot change Runner state, reveals no opponent data, and is
recorded and labeled as model self-report rather than verified Match fact. It
cannot affect scoring except through ordinary provider failures.

\paragraph{Red Phase.} The Orchestrator injects unique ephemeral flags, verifies
their placement, and enables only the declared opponent service paths. Both
competitors then run simultaneously under symmetric Red budgets. Each model may
continue to defend its own service while attacking the opponent's, and may submit
a captured flag to the Orchestrator for verification. A valid submission requires
the correct opponent flag \emph{and} the submitter's own service to pass the
minimum health contract at submission time.

\paragraph{Budgets.} The primary competitive resource is generation consumption,
not elapsed time. Both competitors receive a comparable output-token allowance,
equal maximum turns, equal tool availability, and equal retry and pacing policy;
reasoning tokens are included symmetrically when the provider exposes them.
Wall-clock limits exist only as loose operational backstops. Provider latency is
observed evidence but is not itself a competitive budget, and a transient
provider failure that produces no generation does not consume competitive tokens.

\subsection{Scoring and the best-of-three}

\sandboxer{} uses an outcome-first, lexicographic result rather than an additive
point total. The Orchestrator is the sole authority for flag validity, submission
receipt and verification, service availability, and functional integrity. For a
valid Match, the winner is determined as follows.

\begin{enumerate}
  \item \textbf{Exactly one competitor captures:} that competitor wins, even if
        its own service is no longer healthy at finalization; capture is the
        primary objective and the degraded defense is reported as part of the
        result.
  \item \textbf{Both capture:} if exactly one service is functionally healthy at
        finalization, that competitor wins; if both are healthy, the earlier
        verified submission wins; if neither is healthy, higher measured
        functional availability wins.
  \item \textbf{Neither captures:} higher measured functional availability and
        integrity wins.
  \item \textbf{Unresolved tie:} a fresh, role-swapped tie-break Match is run
        under the same frozen protocol, rather than a coin flip.
\end{enumerate}

Time-to-capture is therefore an explicit ordering criterion only when both
competitors complete equivalent healthy outcomes; it is never converted into an
arbitrary point bonus. Safety and validity outcomes are applied first: an
attributable disqualification awards the Match to the opponent, while an
infrastructure, isolation, or authoritative-telemetry failure invalidates the
Match with no winner.

A public comparison is a \textbf{best-of-three} series, first to two valid Match
wins. Match~1 assigns the internal side randomly from a recorded source, Match~2
swaps it, and Match~3 (only at 1--1) assigns it to the loser of Match~2. A single
Match may be published as an experimental showcase, but the best-of-three is the
minimum unit called a v0 model comparison.

\subsection{Telemetry, evidence, and the Auditor}

\paragraph{Telemetry.} Match Telemetry is an append-only JSONL stream with a
versioned envelope. Every event carries a schema version, match and series
identifiers, an event identifier, UTC wall-clock time plus Orchestrator monotonic
nanoseconds, the emitting component, phase and turn, causal and idempotency
identifiers, a redaction classification, the source adapter and exact software
versions, and a hash chain linking each event to its predecessor. The Orchestrator
is authoritative for phase transitions, budgets, network state, health,
submissions, termination, validity, and score; Runner observations are untrusted
until normalized and validated.

\paragraph{Evidence.} Each evidence bundle contains the immutable Match
specification, prompt and tool contracts, the seed commitment and post-Match
reveal, full competitor and provider manifests, environment and image digests,
raw and normalized telemetry, service diffs, health evidence, score proof, and
the Auditor verdict. Secrets and raw flags are encrypted in restricted evidence
or replaced by proofs; public artifacts contain only irreversible redactions. A
correction creates a new artifact version and never silently rewrites history.

\paragraph{The Match Auditor.} A trusted, independent Auditor supervises every
Match. Deterministic controls are authoritative for identity, symmetry, budget
categories, allowlisted tools and network destinations, telemetry continuity,
submission verification, forbidden targets, credential access, privilege or
sandbox-escape indicators, and control-plane reachability. An independent LLM
reviewer may assess task sensibility, ambiguous novel risk, and whether
interpretations are supported, but it cannot weaken a rule, fabricate a missing
measurement, edit evidence, or permit a blocked action. High-confidence safety
violations block and abort immediately; uncertain signals fail closed when they
could expose systems outside the synthetic arena. Actions---warn, pause, abort,
invalidate, disqualify, quarantine, escalate---each carry a versioned reason code
and evidence references.

\subsection{Competitors, provider, and preflight}

The Competitor is identified by producer and exact model, and the provider is a
declared boundary rather than part of Competitor identity. The first controlled
series runs \texttt{poolside/laguna-s-2.1-free} (Laguna) against
\texttt{meta/muse-spark-1.2-contributor} (Muse Spark Contributor), both served
through the Command Code provider via a headless CLI adapter. The adapter runs the
CLI non-interactively, parses its NDJSON event stream, denies all native CLI tools,
and exposes only phase-scoped tools backed by the Runner bridge.

Before any Match, a \textbf{deterministic provider preflight} records the CLI
binary hash and version, a redacted account fingerprint, a live catalog snapshot,
exact model presence, and the accounting categories actually reported by the
provider, and it live-probes the requested pair to confirm that both models
resolve and are entitled. Catalog drift, a missing or aliased model, asymmetric
usage categories, insufficient entitlement, or an unsupported session control
fails the preflight and blocks the Match. Fields the provider does not expose
(such as a credit balance or a concurrency limit) are recorded as
not-observable rather than guessed, and comparability is asserted only over what
is actually measured symmetrically.

\subsection{Publication pipeline}

All public artifacts are derived views of a frozen, validated evidence bundle,
and every publication step is gated by a human approval. A Result Report
separates \emph{observed} facts, \emph{derived} measurements, \emph{interpreted}
behavior, and \emph{hypotheses}, and states the typing of every analytical claim.
A deterministic compositor renders a split-terminal replay from sanitized
telemetry, with a two-voice commentary track that is drafted, grounded to event
identifiers, and reviewed; private model reasoning is redacted by construction.
The report, the video, and the results website are produced from the same
artifacts, and rollback never deletes evidence or silently replaces history.
Maps, raw traces, flags, and credentials are never exposed publicly.

\section{Limitations}
\label{sec:limitations}

\sandboxer{} is deliberately narrow, and its claims must be read with the
following limitations in mind.

\begin{itemize}
  \item \textbf{Scope of the claim.} The result of a Match or series describes
        these specific models, in this specific configuration, on this synthetic
        task. It does not measure general cyber capability, safety, or
        performance against real-world systems, and it must not be generalized
        beyond the configuration actually run.
  \item \textbf{Synthetic task.} The toy service and the seeded briefs are a
        deliberate simplification. Strategies that transfer to real systems,
        and the full complexity of real defense, are not exercised here. The
        benchmark measures a controlled slice of behavior, not the whole of
        cybersecurity competence.
  \item \textbf{Small samples and variance.} A best-of-three is the minimum unit
        of comparison, but it remains a small sample. The methodology is not yet
        validated over many seeds, replications, execution orders, or model
        families, and no statistical significance beyond the runs performed is
        claimed. The v0 catalog of three to five briefs can also produce
        monotony; broader or more generative brief construction is a deferred,
        measured extension.
  \item \textbf{Provider and scaffold confounders.} Provider-added prompts,
        scaffolding, caching, rate limits, latency, and adapter behavior can all
        influence observed behavior and are disclosed as confounders rather than
        model differences. Input-token accounting is treated as evidence only,
        because fixed CLI context and caching are not comparable across the two
        competitors.
  \item \textbf{Token-budget divergence.} Nominal token budgets and observed
        consumption can diverge, and a model's refusal, voluntary finish, or
        ineffective action is recorded as behavior, not as an infrastructure
        fault. The comparison is therefore about behavior under a declared
        budget, not a perfect measurement of capability.
  \item \textbf{Policy drift.} Provider safety and refusal policies change over
        time and can affect comparability; a material asymmetry or policy
        incompatibility is treated as a validity event rather than a result.
  \item \textbf{Interpretation versus observation.} Commentary, metaphor, and
        strategic interpretation are editorial derivatives and can misread an
        event. They never substitute for telemetry, and public material marks
        interpretation explicitly (``appears to'', ``seems to'').
  \item \textbf{Status of the benchmark.} \sandboxer{} is an experimental
        showcase evaluation. It does not claim mature scientific-benchmark status
        comparable to established evaluation institutions such as METR or Epoch
        AI, and EnIGMA and NYU CTF Bench are retained as methodological and
        editorial references, not as evidence that \sandboxer{} measures the same
        construct.
\end{itemize}

These limitations are not defects to be hidden; they are the boundary of what the
benchmark can responsibly claim. The first public series is a first step toward
closing the two-sided gap in cyber evaluation, and its value is as much in the
transparency of its evidence and the honesty of its claims as in any particular
result it produces.

\newpage
\begin{thebibliography}{9}
\setlength{\itemsep}{2pt}

\bibitem{metr}
T.~Kwa, B.~West, J.~Becker, \emph{et al.},
``Measuring AI ability to complete long tasks,''
arXiv preprint arXiv:2503.14499, 2025.

\bibitem{frontiermath}
E.~Glazer, E.~Erdil, T.~Besiroglu, \emph{et al.},
``FrontierMath: A benchmark for evaluating advanced mathematical reasoning in AI,''
arXiv preprint arXiv:2411.04872, 2024.

\bibitem{inspect}
UK AI Safety Institute,
``Inspect: A framework for large language model evaluations,''
software framework, \url{https://inspect.aisi.org.uk/}.

\bibitem{swebench}
C.~E.~Jimenez, J.~Yang, A.~Wettig, \emph{et al.},
``SWE-bench: Can language models resolve real-world GitHub issues?,''
arXiv preprint arXiv:2310.06770, 2024.

\bibitem{intercode}
J.~Yang, A.~Prabhakar, K.~Narasimhan, and S.~Yao,
``InterCode: Standardizing and benchmarking interactive coding with execution feedback,''
arXiv preprint arXiv:2306.14898, 2023.

\bibitem{cybench}
A.~K.~Zhang, N.~Perry, R.~Dulepet, \emph{et al.},
``Cybench: A framework for evaluating cybersecurity capabilities and risk of language models,''
arXiv preprint arXiv:2408.08926, 2024.

\bibitem{nyuctf}
M.~Shao, S.~Jancheska, M.~Udeshi, \emph{et al.},
``NYU CTF Bench: A scalable open-source benchmark dataset for evaluating LLMs in offensive security,''
arXiv preprint arXiv:2406.05590, 2024.

\bibitem{cybergym}
Z.~Wang, \emph{et al.},
``CyberGym: Evaluating AI agents' real-world cybersecurity capabilities at scale,''
arXiv preprint arXiv:2506.02548, 2025.

\bibitem{enigma}
T.~Abramovich, M.~Udeshi, M.~Shao, \emph{et al.},
``EnIGMA: Interactive tools substantially assist LM agents in finding security vulnerabilities,''
arXiv preprint arXiv:2409.16165, 2024.

\end{thebibliography}

\end{document}
